\chapter{Solutions to Chapter 11: The Central Role of the Propensity Score in Observational Studies for Causal Effects}
\label{sol:chapter11}

\begin{exercisesolution}{Another version of \cref{thm::pscore-dimreduction}}{principal unobserved covariate 的 balancing 性质}
本题的关键是把
\[
W=\{X,Y(1),Y(0)\}
\]
看成一个大的 covariate vector。按照 \cref{def::pscore} 的一般定义，
\[
e(W)=e\{X,Y(1),Y(0)\}=\pr(Z=1\mid W).
\]
我们要证明
\[
Z\ind W\mid e(W),
\]
也就是 \cref{eq::principal-covariate-ind}。由于 $Z$ 是 binary，证明
\[
\pr\{Z=1\mid W,e(W)\}=\pr\{Z=1\mid e(W)\}
\]
即可。

左边很直接。因为 $e(W)$ 是 $W$ 的函数，
\[
\pr\{Z=1\mid W,e(W)\}
=\pr(Z=1\mid W)
=e(W).
\]
右边使用 tower property：
\begin{eqnarray*}
\pr\{Z=1\mid e(W)\}
&=&E\{Z\mid e(W)\}\\
&=&E\{E(Z\mid W,e(W))\mid e(W)\}\\
&=&E\{E(Z\mid W)\mid e(W)\}\\
&=&E\{e(W)\mid e(W)\}\\
&=&e(W).
\end{eqnarray*}
所以两边相等，从而
\[
Z\ind \{Y(1),Y(0),X\}\mid e\{X,Y(1),Y(0)\}.
\]

\paragraph{Remark.}
这个证明没有用到 ignorability。它只是 binary treatment assignment mechanism 的一个概率恒等式：任何变量 $W$ 都可以通过
\[
\pr(Z=1\mid W)
\]
压缩成一个 balancing score。真正有用的地方在于，当 $W=X$ 且 strong ignorability 成立时，这个 balancing score 可由 observed covariates 表示；而本题中的 $W=\{X,Y(1),Y(0)\}$ 包含不可同时观测的 potential outcomes，因此更像一个概念上的基准。
\end{exercisesolution}

\begin{exercisesolution}{Another version of \cref{thm::pscore-dimreduction}}{weak ignorability 的 propensity-score 降维}
固定 $z\in\{0,1\}$。题设只要求
\[
Z\ind Y(z)\mid X,
\]
而不是 joint independence
\[
Z\ind \{Y(1),Y(0)\}\mid X.
\]
因此证明也必须逐个 potential outcome 进行。因为 $Z$ 是 binary，只需证明
\[
\pr\{Z=1\mid Y(z),e(X)\}=\pr\{Z=1\mid e(X)\}.
\]

先看左边。由 tower property，
\begin{eqnarray*}
\pr\{Z=1\mid Y(z),e(X)\}
&=&E\{Z\mid Y(z),e(X)\}\\
&=&E[E\{Z\mid Y(z),e(X),X\}\mid Y(z),e(X)]\\
&=&E[E\{Z\mid Y(z),X\}\mid Y(z),e(X)]\\
&=&E\{E(Z\mid X)\mid Y(z),e(X)\}\\
&=&E\{e(X)\mid Y(z),e(X)\}\\
&=&e(X).
\end{eqnarray*}
其中第四个等号使用了 $Z\ind Y(z)\mid X$。右边同样为
\begin{eqnarray*}
\pr\{Z=1\mid e(X)\}
&=&E\{Z\mid e(X)\}\\
&=&E[E\{Z\mid X,e(X)\}\mid e(X)]\\
&=&E\{E(Z\mid X)\mid e(X)\}\\
&=&E\{e(X)\mid e(X)\}\\
&=&e(X).
\end{eqnarray*}
于是
\[
Z\ind Y(z)\mid e(X).
\]
由于 $z=0,1$ 任意，\cref{thm::pscore-dimension-reduction-weak-ig} 得证。

\paragraph{Remark.}
这个练习对应 \cref{thm::pscore-dimreduction} 的 weak-ignorability 版本。区别在于，weak ignorability 只能逐个 $Y(z)$ 降维；它并不自动给出
\[
Z\ind \{Y(1),Y(0)\}\mid e(X).
\]
这一点和 Chapter 10 中 ``pairwise independence 不推出 joint independence'' 的练习是同一个逻辑警告。
\end{exercisesolution}

\begin{exercisesolution}{More results on the IPW estimators}{HT 的 location non-invariance 与 Hajek 的 normalizing}
先证明 \cref{prop::no-invariance-HT}。记估计 propensity score 为 $\hat e_i=\hat e(X_i)$。HT estimator 为
\[
\hat{\tau}^{\textup{ht}}
=\frac{1}{n}\sumn \frac{Z_iY_i}{\hat e_i}
-\frac{1}{n}\sumn \frac{(1-Z_i)Y_i}{1-\hat e_i}.
\]
若把每个 observed outcome 改为 $Y_i+c$，新的 HT estimator 为
\begin{eqnarray*}
\hat{\tau}^{\textup{ht}}(c)
&=&\frac{1}{n}\sumn \frac{Z_i(Y_i+c)}{\hat e_i}
-\frac{1}{n}\sumn \frac{(1-Z_i)(Y_i+c)}{1-\hat e_i}\\
&=&\hat{\tau}^{\textup{ht}}
+c\left\{
\frac{1}{n}\sumn \frac{Z_i}{\hat e_i}
-\frac{1}{n}\sumn \frac{1-Z_i}{1-\hat e_i}
\right\}\\
&=&\hat{\tau}^{\textup{ht}}
+c(\hat 1_{\textsc{T}}-\hat 1_{\textsc{C}}).
\end{eqnarray*}
这正是 \cref{prop::no-invariance-HT} 的结论。有限样本中一般没有
\[
\hat 1_{\textsc{T}}=\hat 1_{\textsc{C}},
\]
所以 HT estimator 会随着 outcome 的平移而改变。

其次，若使用真实 propensity score $e(X)$，则
\begin{eqnarray*}
E\left\{\frac{1}{n}\sumn \frac{Z_i}{e(X_i)}\right\}
&=&E\left\{\frac{Z}{e(X)}\right\}\\
&=&E\left[E\left\{\frac{Z}{e(X)}\mid X\right\}\right]\\
&=&E\left\{\frac{E(Z\mid X)}{e(X)}\right\}\\
&=&E(1)\\
&=&1.
\end{eqnarray*}
同理，
\begin{eqnarray*}
E\left\{\frac{1}{n}\sumn \frac{1-Z_i}{1-e(X_i)}\right\}
&=&E\left\{\frac{1-Z}{1-e(X)}\right\}\\
&=&E\left[E\left\{\frac{1-Z}{1-e(X)}\mid X\right\}\right]\\
&=&E\left\{\frac{1-E(Z\mid X)}{1-e(X)}\right\}\\
&=&1.
\end{eqnarray*}
这说明 $\hat 1_{\textsc{T}}$ 与 $\hat 1_{\textsc{C}}$ 的目标都是常数 $1$，但 finite-sample realization 不必等于 $1$。

最后证明 Hajek estimator 的 location invariance。写
\[
A_1=\sumn \frac{Z_iY_i}{\hat e_i},\quad
B_1=\sumn \frac{Z_i}{\hat e_i},\qquad
A_0=\sumn \frac{(1-Z_i)Y_i}{1-\hat e_i},\quad
B_0=\sumn \frac{1-Z_i}{1-\hat e_i}.
\]
则
\[
\hat\tau^{\textup{hajek}}=\frac{A_1}{B_1}-\frac{A_0}{B_0}.
\]
若 $Y_i$ 改为 $Y_i+c$，则
\[
A_1\mapsto A_1+cB_1,\qquad A_0\mapsto A_0+cB_0.
\]
因此
\[
\hat\tau^{\textup{hajek}}(c)
=\frac{A_1+cB_1}{B_1}-\frac{A_0+cB_0}{B_0}
=\frac{A_1}{B_1}+c-\frac{A_0}{B_0}-c
=\hat\tau^{\textup{hajek}}.
\]

\paragraph{Remark.}
HT 与 Hajek 的差别不是 identification 层面的差别，而是 finite-sample normalization 的差别。\cref{thm::pscore-weighting} 保证 unnormalized weights 在 expectation 上正确；Hajek estimator 进一步要求样本中 treated 和 control 两侧的总权重都被校准为 $1$。这个小修正常常显著改善有限样本稳定性。
\end{exercisesolution}

\begin{exercisesolution}{Re-analysis of \citet{Rosenbaum::1983JRSSB}}{五个 propensity-score strata 的二元结局分析}
本题的数据已经按估计 propensity score 分成五个 strata。每个 stratum 的总人数都是
\[
303,
\]
所以总体 stratum weight 都是
\[
\pi_k=\frac{303}{1515}=\frac{1}{5}.
\]
令 $\hat p_{k1}$ 和 $\hat p_{k0}$ 分别表示第 $k$ 个 stratum 中 surgical group 与 medical group 的 improvement proportions。基于 propensity-score stratification 的 average causal effect estimator 为
\[
\hat\tau_{\mathrm{strat}}
=\sum_{k=1}^5\pi_k(\hat p_{k1}-\hat p_{k0}).
\]
由 \cref{table::rr1983table1}，
\[
\hat p_{k1}-\hat p_{k0}
=0.19,\ 0.30,\ 0.35,\ 0.41,\ 0.31.
\]
因此
\[
\hat\tau_{\mathrm{strat}}
=\frac{1}{5}(0.19+0.30+0.35+0.41+0.31)
=0.312.
\]
这表示在按 propensity score 分层调整后，surgical therapy 相对于 medical therapy 的 6 个月 functional improvement probability 约高 $31.2$ 个百分点。

接下来给出 standard error。二元 outcome 下，可用
\[
\hat V_{\mathrm{strat}}
=\sum_{k=1}^5\pi_k^2
\left\{
\frac{\hat s_{k1}^2}{n_{k1}}
+\frac{\hat s_{k0}^2}{n_{k0}}
\right\},
\]
其中
\[
\hat s_{kz}^2=\frac{n_{kz}}{n_{kz}-1}\hat p_{kz}(1-\hat p_{kz})
\]
是 binary outcome 的 within-group sample variance。代入表中人数和 proportions，得到
\[
\widehat{\se}(\hat\tau_{\mathrm{strat}})=0.0319.
\]
于是 normal approximation 的 95\% confidence interval 为
\[
0.312\pm 1.96(0.0319)=[0.250,\ 0.374].
\]

下面的代码复现上述计算。由于表中 improvement proportions 只给出两位小数，结果对应 rounded table values。
\begin{lstlisting}
tab <- data.frame(
  stratum = 1:5,
  n1 = c(26, 68, 98, 164, 234),
  p1 = c(0.54, 0.70, 0.70, 0.71, 0.70),
  n0 = c(277, 235, 205, 139, 69),
  p0 = c(0.35, 0.40, 0.35, 0.30, 0.39)
)

tab$nk <- tab$n1 + tab$n0
N <- sum(tab$nk)
tab$pi <- tab$nk / N

est <- sum(tab$pi * (tab$p1 - tab$p0))
s1sq <- tab$n1 / (tab$n1 - 1) * tab$p1 * (1 - tab$p1)
s0sq <- tab$n0 / (tab$n0 - 1) * tab$p0 * (1 - tab$p0)
se <- sqrt(sum(tab$pi^2 * (s1sq / tab$n1 + s0sq / tab$n0)))
c(est = est, se = se, lwr = est - 1.96 * se, upr = est + 1.96 * se)
\end{lstlisting}

\paragraph{Remark.}
这个分析的力量来自表格设计本身：Rosenbaum and Rubin 先用大量 covariates 估计 propensity score，再把问题压缩成五个可比较的 strata。由于原始 individual-level outcomes 不在表中，我们不能恢复未四舍五入的 exact counts；因此这里的数值应理解为基于 published rounded proportions 的近似再分析。
\end{exercisesolution}

\begin{exercisesolution}{Balancing score and propensity score: more theoretical results}{propensity score 是最粗的 balancing score}
我们证明 \cref{thm::balancingscore}。先证明必要性：若 $b(X)$ 是 balancing score，即
\[
Z\ind X\mid b(X),
\]
则 $e(X)$ 是 $b(X)$ 的函数。因为 $b(X)$ 本身是 $X$ 的函数，
\begin{eqnarray*}
e(X)
&=&\pr(Z=1\mid X)\\
&=&\pr\{Z=1\mid X,b(X)\}\\
&=&\pr\{Z=1\mid b(X)\}.
\end{eqnarray*}
右边只依赖于 $b(X)$。因此只要定义
\[
f(u)=\pr\{Z=1\mid b(X)=u\},
\]
就有
\[
e(X)=f\{b(X)\}.
\]
这说明任何 balancing score 都必须至少包含 propensity score 中的信息；也就是说，$b(X)$ 比 $e(X)$ 更细。

再证明充分性：若存在函数 $f(\cdot)$ 使得
\[
e(X)=f\{b(X)\},
\]
则 $b(X)$ 是 balancing score。由于 $Z$ 是 binary，只需证明
\[
\pr\{Z=1\mid X,b(X)\}=\pr\{Z=1\mid b(X)\}.
\]
左边为
\[
\pr\{Z=1\mid X,b(X)\}=\pr(Z=1\mid X)=e(X)=f\{b(X)\}.
\]
右边为
\begin{eqnarray*}
\pr\{Z=1\mid b(X)\}
&=&E\{Z\mid b(X)\}\\
&=&E\{E(Z\mid X)\mid b(X)\}\\
&=&E\{e(X)\mid b(X)\}\\
&=&E[f\{b(X)\}\mid b(X)]\\
&=&f\{b(X)\}.
\end{eqnarray*}
所以两边相等，从而
\[
Z\ind X\mid b(X).
\]
\cref{thm::balancingscore} 得证。

\paragraph{Remark.}
``最粗''的意思是信息压缩的方向：$e(X)$ 可以从任何 balancing score $b(X)$ 中恢复出来。因此 $e(X)$ 是所有 balancing scores 共同必须保留的 assignment information。若为了 subgroup analysis 加上 $X_1$，令
\[
b(X)=\{e(X),X_1\},
\]
那么 $e(X)$ 显然是 $b(X)$ 的函数，所以 $b(X)$ 仍是 balancing score。这就是 \cref{eq::subgroup-ignorability} 背后的概率结构。
\end{exercisesolution}

\begin{exercisesolution}{Some basics of subgroup effects}{subgroup ATE 的 IPW 表示与估计量}
记
\[
G_{x_1}=\mathbb{I}(X_1=x_1),\qquad p_{x_1}=\pr(X_1=x_1).
\]
目标参数为
\[
\tau(x_1)=E\{Y(1)-Y(0)\mid X_1=x_1\}
=\frac{E[G_{x_1}\{Y(1)-Y(0)\}]}{p_{x_1}}.
\]
在 ignorability 和 overlap 条件下，先识别 treated potential-outcome mean 的 subgroup numerator：
\begin{eqnarray*}
E\left\{\frac{G_{x_1}ZY}{e(X)}\right\}
&=&E\left\{\frac{G_{x_1}ZY(1)}{e(X)}\right\}\\
&=&E\left[E\left\{\frac{G_{x_1}ZY(1)}{e(X)}\mid X\right\}\right]\\
&=&E\left[\frac{G_{x_1}}{e(X)}E\{ZY(1)\mid X\}\right]\\
&=&E\left[\frac{G_{x_1}}{e(X)}E(Z\mid X)E\{Y(1)\mid X\}\right]\\
&=&E[G_{x_1}E\{Y(1)\mid X\}]\\
&=&E\{G_{x_1}Y(1)\}.
\end{eqnarray*}
同理，
\[
E\left\{\frac{G_{x_1}(1-Z)Y}{1-e(X)}\right\}
=E\{G_{x_1}Y(0)\}.
\]
两式相减并除以 $p_{x_1}$，得到题目所给的 identification formula：
\[
\tau(x_1)
=E\left\{
\frac{\mathbb{I}(X_1=x_1)ZY}{e(X)}
-\frac{\mathbb{I}(X_1=x_1)(1-Z)Y}{1-e(X)}
\right\}\Big/\pr(X_1=x_1).
\]

令 $\hat e_i=\hat e(X_i)$，$G_i(x_1)=\mathbb{I}(X_{i1}=x_1)$，以及
\[
n_{x_1}=\sumn G_i(x_1).
\]
对应的 subgroup HT estimator 可以写成
\[
\hat\tau^{\textup{ht}}(x_1)
=\frac{1}{n_{x_1}}\sumn G_i(x_1)
\left\{
\frac{Z_iY_i}{\hat e_i}
-\frac{(1-Z_i)Y_i}{1-\hat e_i}
\right\}.
\]
等价地，它是用 empirical subgroup probability $n_{x_1}/n$ 估计 $p_{x_1}$ 后的 plug-in estimator。

对应的 subgroup Hajek estimator 分别 normalizes treated 和 control weights：
\[
\hat\tau^{\textup{hajek}}(x_1)
=
\frac{\sumn G_i(x_1)Z_iY_i/\hat e_i}
{\sumn G_i(x_1)Z_i/\hat e_i}
-
\frac{\sumn G_i(x_1)(1-Z_i)Y_i/(1-\hat e_i)}
{\sumn G_i(x_1)(1-Z_i)/(1-\hat e_i)}.
\]
它把 subgroup 内的 treated pseudo-population 和 control pseudo-population 分别校准到同一个 subgroup。

\paragraph{Remark.}
subgroup analysis 的一个常见误区是只在 propensity score model 里加入 subgroup indicator，却没有在估计阶段明确 target population。上面的公式说明，若目标是 $\tau(x_1)$，则 weighting 不仅要调整 $X$，还要把 averaging distribution 限制在 $X_1=x_1$ 的 subgroup 内。换言之，propensity score 解决 comparability，indicator $G_{x_1}$ 决定目标人群。
\end{exercisesolution}
